\documentclass[a4paper,11pt]{article}
\usepackage[utf8]{inputenc}
\usepackage[T1]{fontenc}
\usepackage[french]{babel}
\usepackage[left=2cm,right=2cm,top=2cm,bottom=2cm]{geometry}
\usepackage{xcolor}
\usepackage{hyperref}
\usepackage{fontawesome5}
\usepackage{parskip}

% --- Couleurs Sopra Steria ---
\definecolor{primary}{RGB}{190, 30, 45}
\hypersetup{colorlinks=true, urlcolor=primary}

\begin{document}

% --- EXPÉDITEUR ---
\noindent
\begin{minipage}[t]{0.5\textwidth}
    \textbf{Loïc Aron MBASSI EWOLO}\\
    Élève-Ingénieur Bac+5 -- Double diplôme ENSEA\\[3pt]
    \faMapMarker\ Cergy, France\\
    \faPhone\ (+33) 7 59 52 33 98\\
    \faEnvelope\ \href{mailto:loicaron.mbassiewolo@ensea.fr}{loicaron.mbassiewolo@ensea.fr}
\end{minipage}
\hfill
% --- DATE ---
\begin{minipage}[t]{0.4\textwidth}
    \raggedleft
    Cergy, le 2 mars 2026
\end{minipage}

\vspace{1.2cm}

% --- DESTINATAIRE ---
\hfill
\begin{minipage}[t]{0.5\textwidth}
    \raggedleft
    \textbf{Sopra Steria}\\
    Aéroline Air Traffic Management\\
    Colomiers, France
\end{minipage}

\vspace{1cm}

\noindent\textbf{Objet :} Candidature en alternance -- Développeur Full Stack IA

\vspace{0.8cm}

Madame, Monsieur,

Concevoir un système multi-agents avec RAG et LangChain, le déployer via FastAPI et Docker, puis le présenter à une équipe de recherche : c'est le projet que j'ai mené ces derniers mois en parallèle de mon cursus d'ingénieur. Cette expérience m'a préparé aux missions que vous proposez au sein d'Aéroline, où l'IA générative est mise au service du contrôle aérien.

En double diplôme à l'ENSEA après un cursus en Mathématiques Appliquées et Génie Logiciel à l'École Polytechnique de Yaoundé, je travaille sur les technologies LLM/RAG depuis 2024. Mon projet le plus abouti est un système multi-agents agricole : 4 agents IA collaboratifs orchestrés via Google ADK, alimentés par RAG pour l'enrichissement contextuel, exposés par une API FastAPI et déployés avec Docker. L'architecture, la résolution de conflits inter-agents et l'intégration d'APIs externes (météo, marchés) m'ont confronté aux mêmes enjeux que vos agents support aux opérateurs du contrôle aérien : fiabilité des réponses, pertinence contextuelle et industrialisation du Proof-of-Concept.

Sur le volet NLP et Transformers, j'ai fine-tuné des modèles BERT et GPT-2 pour la détection de contenus sexistes (compétition MLPC), et développé une application de reconnaissance vocale sur une langue peu dotée en exploitant des modèles pré-entraînés HuggingFace (projet YembaAudio). Mon stage au MILA (Institut Québécois d'IA) m'a formé à la rigueur de la recherche appliquée en IA : reproduction d'expériences SOTA en PyTorch, analyse mathématique et rédaction de rapports de synthèse. En parallèle, +3 ans d'expérience Full Stack (plateforme de télémédecine en production, outil open source de génération de code avec intégration LLM) garantissent ma capacité à livrer du code industrialisable.

La perspective de travailler sur des problématiques concrètes de gestion du trafic aérien, dans un cadre Agile avec une chaîne DevOps, correspond à la suite logique de mon parcours. Je suis disponible dès septembre 2026 pour une alternance et mobile pour rejoindre le site de Colomiers/Toulouse. La possibilité d'embauche en CDI renforce ma motivation à m'investir pleinement dans cette mission.

Je reste à votre disposition pour un entretien afin de discuter de ma candidature.

Je vous prie d'agréer, Madame, Monsieur, l'expression de mes salutations distinguées.

\vspace{0.8cm}

\textbf{Loïc Aron MBASSI EWOLO}\\
\textit{Élève-Ingénieur ENSEA / Polytechnique Yaoundé}

\end{document}
